% !TEX root = ../main.tex
% !TEX program = pdflatex
\subsection{Определение размера строки кэша}

Строка (линия) кэша~--- минимальная порция данных, передаваемая между уровнями памяти. Размер строки определялся двумя независимыми методами.

\textbf{Метод А: обход с шагом.} Массив размером 256~МБ (заведомо больше L2) обходится последовательно с шагом $s$ байт, к каждому затронутому элементу прибавляется единица: \code{for (i = 0; i < n; i += s/4) arr[i]++}. Измеряется время одного полного прохода. Пока шаг меньше размера строки $L$, затрагивается каждая строка массива, из памяти загружается весь массив, и время прохода почти не зависит от шага. Когда $s \ge L$, часть строк пропускается, и каждое удвоение шага вдвое уменьшает объём загружаемых данных и время.

\fig{line_a.png}{Метод А: время прохода массива 256~МБ в зависимости от шага}{fig:line_a}

Как видно на рисунке~\ref{fig:line_a}, при шагах 4 и 8~байт время ограничено числом выполняемых команд (каждое удвоение шага вдвое уменьшает число итераций). Далее при шагах от 16 до 128~байт время остаётся почти постоянным (7,1--8,8~мс)~--- пропускная способность памяти тратится на загрузку одних и тех же строк. Начиная с шага 256~байт время падает вдвое на каждое удвоение шага: 4,8; 2,6; 1,2; 0,6~мс. Значит, строки, загружаемые из памяти, имеют размер 128~байт.

\textbf{Метод Б: два чтения в блоке.} Массив разбивается на блоки по 4~КБ, порядок блоков перемешивается случайно. В цепочку включаются по два элемента каждого блока: с нулевым смещением и со смещением $d$. Если $d < L$, оба элемента лежат в одной строке и второе обращение попадает в кэш; если $d \ge L$, на каждый блок приходится два промаха. Измеряется время на один блок. Чтобы определить строку разных уровней, эксперимент проводился для массива 256~МБ (промахи идут в оперативную память) и для массива 8~МБ (данные в L2, промахи только в L1).

\fig{line_l2_b.png}{Метод Б, массив 8~МБ (данные в L2): время на блок в зависимости от расстояния $d$}{fig:line_l2}

\fig{line_b.png}{Метод Б, массив 256~МБ (данные в памяти): время на блок в зависимости от расстояния $d$}{fig:line_b}

Для массива 8~МБ (рисунок~\ref{fig:line_l2}) при $d \le 32$~байт время на блок около 8~нс (один промах L1 и одно попадание), а уже при $d = 64$~байт оно возрастает до 12~нс и дальше остаётся на уровне 11--14~нс (два промаха L1). Следовательно, \textbf{размер строки кэша L1 равен 64~байтам}.

Для массива 256~МБ (рисунок~\ref{fig:line_b}) при $d \le 32$~байт время на блок около 110~нс (одно обращение к памяти), при $d = 64$ и $128$~байт время растёт частично (124 и 155~нс), а при $d = 256$~байт достигает 309~нс~--- двух полных обращений к памяти. Промежуточные значения объясняются тем, что при $d = 64$ второй элемент приходится на вторую половину 128-байтной строки L2, которая поступает чуть позже первой, а при $d = 128$ соседнюю строку частично подгружает аппаратная предвыборка. Снижение времени при $d \ge 512$ связано с тем, что оба обращения попадают в одну и ту же открытую строку микросхемы DRAM. Вместе с методом А это даёт \textbf{размер строки кэша L2 128~байт}, что совпадает со значением \code{hw.cachelinesize = 128}, которое сообщает ОС.

Итак, строка кэша L1 составляет 64~байта (16 элементов \code{int}), а обмен между L2 и памятью идёт строками по 128~байт (32 элемента \code{int}).
